\section{The original completed point field gives the genuine Laurent Gauss product point}

\begin{lemma}\label{lem:native-point-completed-field-coefficients}
\leanok
\lean{EnvelopeProperty.Berkovich.SchemePointLaurentGauss.PointField}
\lean{EnvelopeProperty.Berkovich.SchemePointLaurentGauss.coefficientMap}
\lean{EnvelopeProperty.Berkovich.SchemePointLaurentGauss.coefficientMap_spec}
\lean{EnvelopeProperty.Berkovich.SchemePointLaurentGauss.coefficientMap_val}
Let $x$ be any actual relative valued point of an original scheme $X$
over a commutative valued ring $(R,v)$. Its actual completed stalk residue
field $H(x)$ has a canonical coefficient map $\varphi_x:R\to H(x)$.
The entire scheme coefficient triangle commutes and
$|\varphi_x(a)|=v(a)$ for every original $a\in R$.
\uses{scheme:actual-native-completed-field-point}
\end{lemma}
\begin{proof}
Compose the original canonical field-valued point
$\operatorname{Spec}H(x)\to X$ with the original structure map and use the
full faithfulness of the affine spectrum functor to obtain the actual
ring map. Mapping the actual norm point of $H(x)$ through its field-valued
point recovers the entire original point $x$. Its prescribed-base equation
therefore proves equality of the two original base points on
$\operatorname{Spec}R$. Evaluate that whole equality on every original
coefficient to get the asserted norm formula.
\end{proof}

\begin{definition}\label{def:native-point-laurent-gauss-family}
\leanok
\lean{EnvelopeProperty.Berkovich.SchemePointLaurentGauss.point}
\lean{EnvelopeProperty.Berkovich.SchemePointLaurentGauss.originalProjection}
\lean{EnvelopeProperty.Berkovich.SchemePointLaurentGauss.parameterProjection}
\lean{EnvelopeProperty.Berkovich.SchemePointLaurentGauss.ordinaryParameterProjection}
\lean{EnvelopeProperty.Berkovich.SchemePointLaurentGauss.sourceOver}
\lean{EnvelopeProperty.Berkovich.SchemePointLaurentGauss.originalOverMap}
\lean{EnvelopeProperty.Berkovich.SchemePointLaurentGauss.parameterOverMap}
\lean{EnvelopeProperty.Berkovich.SchemePointLaurentGauss.productMap}
\lean{EnvelopeProperty.Berkovich.SchemePointLaurentGauss.relativeSourcePoint}
\lean{EnvelopeProperty.Berkovich.SchemePointLaurentGauss.productPoint}
\lean{EnvelopeProperty.Berkovich.SchemePointLaurentGauss.productMap_fst}
\lean{EnvelopeProperty.Berkovich.SchemePointLaurentGauss.productMap_snd}
For $r>0$, take the actual Laurent Gauss point $z$ of
$\operatorname{Spec}H(x)[T,T^{-1}]$ at radius $r$. Its original projection
to $X$ is the spectrum of the actual coefficient inclusion followed by
the original completed field point. Its Laurent parameter projection is
the spectrum of the actual Laurent extension of $\varphi_x$.
\uses{lem:native-point-completed-field-coefficients,thm:laurent-gauss-coefficient-naturality,spectrum:actual-laurent-gauss-valuation,def:native-multiplicative-group}
\end{definition}
\begin{proof}
Both maps are actual scheme maps from the same Laurent spectrum. Their
composites to the original $R$-spectrum agree, since their ring coefficient
composites are the same $a\mapsto C(\varphi_x(a))$. They therefore give
the original product map over $R$, using the actual Cartesian product.
No valued point or projection matching is supplied as an input.
\end{proof}

\begin{theorem}\label{thm:native-point-laurent-gauss-projections}
\leanok
\lean{EnvelopeProperty.Berkovich.SchemePointLaurentGauss.originalProjection_point}
\lean{EnvelopeProperty.Berkovich.SchemePointLaurentGauss.parameterProjection_point}
\lean{EnvelopeProperty.Berkovich.SchemePointLaurentGauss.ordinaryParameterProjection_point}
\lean{EnvelopeProperty.Berkovich.SchemePointLaurentGauss.productPoint_fst}
\lean{EnvelopeProperty.Berkovich.SchemePointLaurentGauss.productPoint_snd}
The whole original projection of $z$ to $X$ equals $x$. Its whole Laurent
parameter projection is the actual Laurent Gauss point of $v$ at radius
$r$. Consequently its projection to the ordinary parameter line is the
actual polynomial Gauss point of $v$ at radius $r$.
\uses{def:native-point-laurent-gauss-family,lem:native-point-completed-field-coefficients,thm:laurent-gauss-coefficient-naturality}
\end{theorem}
\begin{proof}
Actual Laurent Gauss restriction on the coefficient field equals its
original norm valuation. The canonical completed field point therefore
recovers the entire original native $x$. The original Laurent coefficient
map preserves every coefficient valuation by the preceding lemma, so
whole Laurent Gauss naturality proves the second equality. Restricting it
through the original polynomial localization proves the third equality.
All equalities concern the entire valued points and all original
coefficients, not only their scheme supports or parameter value.
\end{proof}

\begin{theorem}\label{thm:native-point-laurent-gauss-section-values}
\leanok
\lean{EnvelopeProperty.Berkovich.SchemePointLaurentGauss.section_laurent_value}
Let $U$ be any original open containing the support of $x$, and let
$f=\sum_i a_iT^i$ be any Laurent polynomial with original coefficients
$a_i\in\Gamma(X,U)$. Evaluating every coefficient through the actual
original section map into $H(x)$, the actual Gauss value is
$\max_{i\in\operatorname{supp}f}|a_i(x)|r^i$.
\uses{def:native-point-laurent-gauss-family,cor:laurent-gauss-original-coefficient-maximum}
\end{theorem}
\begin{proof}
The original section-to-completed-field map retains the actual native
value of every original section. Apply the proved whole Laurent
coefficient maximum formula to this actual ring map. The support is the
original support of $f$, so zero values of some coefficients are retained
without a nonvanishing premise.
\end{proof}
